\documentclass[11pt,letterpaper]{article}
\usepackage{hanzo-defense}

\doctitle{Confidential Computing by Default:\\An FHE-Native Chain and\\Agent Runtime}
\docid{HAI-DEF-2026-014}
\docversion{Version 1.0}
\docdate{2026}
\docauthors{Zach Kelling --- Hanzo Team}
\docorgs{Hanzo Industries\\research@hanzo.ai}

\begin{document}
\makecover

\begin{abstract}
Every system that processes sensitive data has a moment where it holds the data
in the clear---and that moment is the insider risk, the cross-domain leak, and the
blast radius of every breach. Fully homomorphic encryption (FHE) removes that
moment: computation runs directly on ciphertext, and only a threshold of keyholders
can decrypt the result. This paper describes Hanzo's design for making FHE the
\emph{default} mode of computation across the two layers that normally demand
plaintext---a blockchain's state and rules, and an autonomous agent's actions---so
that state stays encrypted on-chain, policies are evaluated over ciphertext, and
agents act on data they are never able to read. We give the cryptographic basis
(BFV/BGV/CKKS/TFHE with threshold decryption), the GPU acceleration that makes it
practical, the post-quantum anchoring that protects it for the long term, the
defense applications it unlocks, and an explicit, honest assessment of which
components are production today versus an active research frontier.
\end{abstract}

\paragraph{Executive Summary.}
Normally, to compute on data you must first decrypt it---so somewhere in every
system there is a component holding the secret in the clear, and that component is
what an adversary, an insider, or a subpoena goes after. FHE lets the computer do
the work \emph{while the data stays encrypted}: it takes in ciphertext, produces a
ciphertext answer, and never sees the plaintext. Hanzo pushes this idea through the
whole stack. The blockchain stores encrypted state and can enforce rules on
encrypted data; agents can act on encrypted data and return encrypted answers; and
only a quorum of keyholders can ever decrypt---no single party holds the key. For
defense this is directly useful: allies can compute a joint answer without sharing
raw data, a cross-domain guard can decide whether something is releasable without
seeing its content, and sensitive workloads can run on infrastructure you do not
fully control without exposing the data to the host. FHE's historical objection is
speed; the practical answer here is GPU acceleration plus running only the
sensitive part of a computation under encryption. The result is a stack where
confidentiality is the default state of data in use, not just data at rest or in
transit.

\tableofcontents
\newpage

% ============================================================================
\section{The Plaintext Gap}
% ============================================================================

We protect data in two of its three states well and the third state badly. Data
\emph{at rest} is encrypted on disk; data \emph{in transit} is encrypted on the
wire. But data \emph{in use}---the moment a program actually computes on it---is
almost always plaintext in memory. Every database query that filters records,
every policy engine that decides access, every model that runs inference, every
smart contract that checks a balance, does so on cleartext. That cleartext moment
is the structural weakness behind a large share of real incidents: the compromised
process, the malicious insider, the memory-scraping implant, the cross-domain
guard that has to read the data to rule on it, the cloud host that can see its
tenants' workloads.

Confidential computing has tried to shrink this gap with hardware enclaves (trusted
execution environments). Enclaves help, but they relocate trust rather than remove
it: you must trust the silicon vendor, the attestation chain, and the absence of
side channels, and the data is still plaintext \emph{inside} the enclave. Fully
homomorphic encryption attacks the gap directly. With FHE the data is never
decrypted to be computed on at all---not in the enclave, not in memory, nowhere.
The computation is performed on ciphertext and yields a ciphertext result. The
plaintext gap closes.

% ============================================================================
\section{FHE in One Page}
% ============================================================================

Fully homomorphic encryption is an encryption scheme with an extra property: there
exist operations on ciphertexts that correspond to operations on the underlying
plaintexts. Concretely, given encryptions of $a$ and $b$, anyone---without the
key---can compute an encryption of $a+b$ and of $a\times b$. Addition and
multiplication are a complete basis, so in principle any computable function can be
evaluated homomorphically by expressing it as an arithmetic (or boolean) circuit.
The party doing the computation learns nothing; only a holder of the secret key can
decrypt the output.

\paragraph{The scheme families and what each is for.}
\begin{itemize}
  \item \textbf{BFV / BGV} --- exact integer arithmetic. Right for counts, sums,
  comparisons, exact matching, and policy predicates over integers.
  \item \textbf{CKKS} --- approximate arithmetic over real/complex numbers, with
  SIMD ``batching'' that packs thousands of values into one ciphertext. Right for
  statistics, signal processing, and machine-learning inference.
  \item \textbf{TFHE} --- fast boolean gates and programmable bootstrapping. Right
  for branching, comparisons, and arbitrary logic, including data-dependent control
  flow.
\end{itemize}

\paragraph{Bootstrapping.} Each homomorphic operation adds noise; after enough
operations the noise would corrupt the result. \emph{Bootstrapping} refreshes a
ciphertext to a low-noise state and is what makes encryption ``fully''
homomorphic---it permits unbounded computation depth. Bootstrapping is the
expensive step, and it is the principal target of the acceleration discussed in
\S\ref{sec:speed}.

\paragraph{Threshold decryption.} The secret key need not exist in one place. With
\emph{threshold} FHE the key is split across $n$ parties so that any $t$ of them can
jointly decrypt but no coalition smaller than $t$ can. Combined with FHE this gives
a system in which \emph{no single party ever holds both the data and the ability to
reveal it}---the computation runs blind, and revealing the answer requires a quorum.
This property is what makes the design defensible for multi-party and adversarial
settings.

% ============================================================================
\section{Confidential State: An FHE-Native Chain}
% ============================================================================

A blockchain is normally the \emph{opposite} of confidential: every validator
re-executes every transaction against shared plaintext state, which is precisely
why public chains leak everything. An FHE-native chain inverts this. State is
stored as ciphertext; transactions carry encrypted inputs; and the state-transition
rules are evaluated homomorphically, so validators advance the ledger
\emph{without ever seeing the values they are agreeing on}.

\begin{itemize}
  \item \textbf{Encrypted state.} Balances, ownership records, order books,
  classifications---whatever the application holds---live on-chain as ciphertext.
  The authoritative record exists, is immutable and auditable, and is unreadable to
  anyone without a decryption quorum.
  \item \textbf{Rules on ciphertext.} A transfer rule such as ``permit only if the
  sender's balance covers the amount'' is a comparison; under BFV/TFHE that
  comparison is evaluated on the encrypted balance and encrypted amount, yielding an
  encrypted yes/no that gates the state update. The rule is enforced without the
  rule-evaluator learning the balance.
  \item \textbf{Threshold-gated disclosure.} When a value must be revealed---a
  settlement, an audit, a court order, a coalition release decision---a threshold of
  keyholders cooperates to decrypt exactly that value, and the decryption itself is
  recorded. Disclosure becomes an explicit, quorum-authorized, logged event rather
  than a default.
  \item \textbf{Post-quantum anchoring.} The chain's signatures and transport are
  post-quantum (ML-KEM/ML-DSA), so the encrypted record and its provenance survive a
  future quantum adversary---essential because ciphertext recorded immutably today
  is exactly the ``harvest now, decrypt later'' target (see the full-spectrum cyber
  and ZAP papers).
\end{itemize}

The chain thus provides a tamper-evident, post-quantum, immutable audit trail
\emph{of encrypted facts}: you can prove what happened and in what order without
ever exposing what the facts were, until a quorum decides otherwise.

% ============================================================================
\section{Confidential Action: FHE-Native Agents}
% ============================================================================

The second layer that normally demands plaintext is automation. An autonomous
agent that triages intelligence, screens transactions, routes logistics, or
enforces a policy has to ``see'' the data to act on it---and so the agent runtime
becomes another plaintext concentration point, now with the additional risk that
the model itself might memorize or exfiltrate. FHE breaks this too.

An FHE-native agent operates on encrypted inputs and emits encrypted outputs. It
can evaluate the encrypted state of the chain, apply encrypted rules, and produce
an encrypted decision or an encrypted action---all without the agent operator (or
the model host) ever seeing the plaintext. Where the agent must run a learned model
rather than a fixed rule, CKKS-based encrypted inference lets a (suitably
structured) model classify or score encrypted inputs and return an encrypted
result. The agent is then \emph{operator-blind}: it does useful work on data it
cannot read, on infrastructure that cannot read it either.

This composes with Hanzo's broader agentic runtime (the agentic-AI and edge papers):
the same orchestration of specialized agents and tool calls runs, but the data
plane underneath it is ciphertext. Autonomy and confidentiality stop being a
trade-off.

% ============================================================================
\section{Programmable Confidentiality: Rules and Policies on Ciphertext}
% ============================================================================

The unifying idea is that a \emph{rule} is just a function, and any function can in
principle be evaluated under FHE. So the policies that govern a system---not only
its data---can run on ciphertext. This is the capability the rest of the paper
builds on, because most high-consequence decisions are policy decisions over
sensitive inputs:

\begin{itemize}
  \item \textbf{Releasability} --- ``may this datum cross to that security domain /
  that coalition partner?'' decided on the encrypted datum (see cross-domain paper).
  \item \textbf{Eligibility / authorization} --- ``does this party satisfy the
  conditions?'' decided without revealing the party's attributes.
  \item \textbf{Thresholds and limits} --- ``is the aggregate over/under the line?''
  decided without revealing the individual contributions.
  \item \textbf{Matching} --- ``do these two sealed values correspond?'' (sealed-bid
  auctions, private set intersection, deconfliction) decided without opening either
  side.
\end{itemize}

In each case the \emph{decision} is computed and enforced while the \emph{inputs}
remain encrypted, and only the (small, controlled) decision output is ever a
candidate for threshold decryption. Confidentiality becomes programmable: you write
the rule once and it runs blind.

% ============================================================================
\section{Defense Applications}
% ============================================================================

The combination of an encrypted ledger and operator-blind agents answers a set of
defense problems that have no clean classical solution:

\begin{itemize}
  \item \textbf{Coalition analytics without disclosure.} Partners compute a joint
  result---shared targets, overlapping indicators, combined logistics---over their
  combined data while each party's raw data stays encrypted. Nobody surrenders
  sources or methods to get the joint answer, and the chain records that the
  computation occurred.
  \item \textbf{Cross-domain transfer on ciphertext.} A guard evaluates a release
  policy against encrypted content and decides movement between security levels
  without ever seeing the content, removing the guard itself as a data tap.
  \item \textbf{Private intelligence fusion.} Multi-source correlation runs while
  the sources stay encrypted, so a breach of the fusion system does not expose the
  underlying collection.
  \item \textbf{Encrypted ML inference.} A model classifies encrypted imagery or
  signals: the data owner never exposes the input and the model operator never sees
  it---useful when the data is sensitive, the model is sensitive, or both.
  \item \textbf{Encrypted command, targeting, and logistics state.} Orders,
  target lists, and movement tables live as encrypted on-chain state with an
  immutable, post-quantum audit trail; agents act on them under encryption; reveal
  is quorum-gated and logged.
  \item \textbf{Computation on untrusted infrastructure.} Sensitive workloads run on
  contractor, coalition, or commercial-cloud hardware without exposing the data to
  the host---widening where a mission can compute.
\end{itemize}

% ============================================================================
\section{Why the Speed Objection No Longer Holds}
\label{sec:speed}
% ============================================================================

FHE's reputation for impracticality is a decade out of date, and three things close
the gap in this stack:

\begin{enumerate}
  \item \textbf{GPU acceleration.} The dominant cost in FHE is the
  number-theoretic transform (NTT) at the heart of polynomial multiplication and
  bootstrapping. These are massively parallel and map naturally onto GPUs; the same
  native-GPU foundation that accelerates the chain's cryptography accelerates the
  FHE layer, turning bootstrapping from seconds toward the practical range.
  \item \textbf{Batching.} CKKS and BGV pack thousands of plaintext slots into one
  ciphertext, so a single homomorphic operation does thousands of useful operations
  at once---ideal for the vector and tensor workloads in analytics and inference.
  \item \textbf{Selective FHE.} Most computations have a small sensitive core inside
  a large non-sensitive shell. You do not run the whole program under FHE---only the
  predicate or the field that must stay private. The expensive primitive is spent
  precisely where it buys confidentiality and nowhere else.
\end{enumerate}

The honest framing is not ``FHE is now as fast as plaintext''---it is not---but
``the workloads that matter for confidential rules, matching, gating, and bounded
inference are now practical, and GPU acceleration keeps widening the set.''

% ============================================================================
\section{Key Custody Without a Single Point of Trust}
% ============================================================================

A confidentiality system is only as strong as its key handling. The design holds
the decryption key nowhere: threshold decryption splits it $t$-of-$n$ across
independent parties (validators, agencies, coalition members) so that decrypting
any value requires an explicit quorum and no minority---and no single compromised
host, insider, or seized server---can reveal anything. This is the structural
answer to ``who can see the data?'': by construction, \emph{nobody}, until enough
authorized parties agree, and that agreement is itself recorded on the chain.

% ============================================================================
\section{Composition With the Rest of the Stack}
% ============================================================================

FHE here is not a bolt-on library; it is one primitive in a stack Hanzo owns
end-to-end, which is what makes the ``native'' claim real:

\begin{itemize}
  \item \textbf{Native-GPU chain} provides the encrypted, immutable, post-quantum
  ledger and the GPU substrate the FHE layer rides on.
  \item \textbf{Post-quantum transport and identity} (ZAP, ML-KEM/ML-DSA) protect
  the ciphertext in motion and bind it to verifiable identities.
  \item \textbf{Agentic runtime} supplies the operator-blind automation that acts on
  the encrypted state.
  \item \textbf{Owned models} (the Zen family) supply the inference that, under
  CKKS, can run on encrypted inputs.
\end{itemize}

Because all four are owned and open, the confidentiality guarantee is auditable
rather than asserted: a defense customer can inspect the whole path the data takes
and confirm there is no point where it is silently in the clear.

% ============================================================================
\section{Honest Maturity Assessment}
% ============================================================================

This is a capability paper, and the credible version of it is explicit about what
runs today versus what is an active frontier. No overstatement:

\begin{itemize}
  \item \textbf{Production-grade today:} the GPU-accelerated FHE primitives
  (BFV/BGV/CKKS/TFHE), threshold decryption, the native-GPU post-quantum chain, and
  the agentic runtime each exist and run. Bounded-depth confidential workloads
  ---rule and policy predicates, comparisons, matching/deconfliction, aggregate
  thresholds, and inference on suitably structured models---are practical now.
  \item \textbf{Active frontier:} fully general, arbitrary-depth FHE-native
  programmability at plaintext-class throughput is not solved---by anyone---and we do
  not claim it. The performance envelope is widening (faster bootstrapping, better
  GPU kernels, hybrid MPC/FHE designs), and the architecture is built so that
  workloads move from ``frontier'' to ``practical'' as that envelope grows, without
  redesign.
\end{itemize}

The strategic point stands regardless of where the line currently sits: Hanzo owns
every layer the line runs through, so as FHE performance improves the gains land in
a stack a customer already controls---rather than waiting on a closed vendor to
expose the capability.

% ============================================================================
\section{Conclusion}
% ============================================================================

Confidentiality of data \emph{in use} is the unsolved third of the problem, and it
is the one that matters most for coalition operations, cross-domain movement,
intelligence fusion, and computing on infrastructure you do not own. FHE closes the
plaintext gap by computing on ciphertext; threshold decryption removes the single
point of trust; GPU acceleration and selective application make the important
workloads practical; and post-quantum anchoring protects the encrypted record for
the long term. Hanzo's contribution is to make this the default across a chain and
an agent runtime it owns end-to-end, so that ``private ciphertext for every
operation'' is an architectural property of the stack rather than a feature of one
application. Compute on data you are not allowed to see---and let the rules,
the ledger, and the agents all run blind.

\end{document}
